% 来源：成文/A\_08\_模型建立与求解.md ＋ 成文/A\_09\_结果分析.md（原 5.5，并入本章）
\section{模型建立与求解}
% 本章子节编号自 5.0 起（计数器置 -1：首个子节即 5.0）—— 属**计数器机制**，非写死编号
\setcounter{subsection}{-1}
\subsection{总体模型与求解框架}
四问共用一个物理过程：圆柱形药材在热风环境中，内部以导热方式升温、以扩散方式失水，两者通过物性相互影响。本文据此建立一维轴对称圆柱非稳态传热传质耦合模型，其总体结构如图~\ref{fig:5-1} 所示。


% 图5-1 ＝ F-01（F 号）／图01（PDF 文件索引号）／建模总体流程图
\begin{figure}[htbp]
  \centering
\includegraphics[width=1.000\textwidth]{../../05\_成品图/图01-总体建模流程图.pdf}
  \caption{建模总体流程图}\label{fig:5-1}
\end{figure}

模型由“题面与附件数据 → 守恒律建立控制方程 → 元体平衡离散 → 全隐式求解 → 检验与灵敏度”五段构成，四问沿同一主干分叉。由图~\ref{fig:5-1} 可见，四问的差别不体现在方程组形式上，而只体现在四组旋钮的位置：物性体系、耦合方式、时程与终止判据、几何是否收缩。本文是“一个模型、四组旋钮”的闭环结构，而非四个独立模型。

\subsubsection{控制方程族（四问共用形式）}
在圆柱坐标系下，令 $T(r,t)$ 为温度、$C(r,t)$ 为干基含水率，取散度形式（保留曲率项 $1/r$）：

\begin{equation}
  \rho c_p\frac{\partial T}{\partial t}=\frac{1}{r}\frac{\partial}{\partial r}\Big(k\,r\frac{\partial T}{\partial r}\Big),\qquad \frac{\partial C}{\partial t}=\frac{1}{r}\frac{\partial}{\partial r}\Big(D\,r\frac{\partial C}{\partial r}\Big)
  \label{eq:5-1}
\end{equation}
其中 $\rho,c_p,k$ 与 $D$ 的取值按问分套：问题一取常值，问题二、三取附录 3 的经验式，问题四取附录 4 的经验式并叠加几何收缩。需要强调的是，导热系数 $k$ 必须位于散度算子之内；写成“$k$ 乘拉普拉斯算子”的形式在变物性下不再等价，会丢掉由物性梯度产生的附加通量。

\subsubsection{求解与验证路线}
空间上统一采用元体平衡（单元中心有限体积）离散，其选用依据有二：相邻控制体在公共界面使用同一数值通量，求和时内部通量两两抵消，区域总量的变化只取决于最外层边界通量，故守恒性在代数层面严格成立且与网格是否均匀无关；该格式导出的系数矩阵为 M 矩阵，离散极值原理随之成立，故固定计算域（问题一至问题三）下温度不越界、含水率沿径向单调属理论必然，而非经验观察（问题四的移动边界与压缩温升另见 5.4.5）。

上述“内部通量两两抵消”可写成对任一控制体 $i$ 与任一守恒量 $\phi$（取温度或干基含水率）都成立的通量平衡式：

\begin{equation}
  \frac{V_i}{\Delta t}\big(\phi_i^{\,n+1}-\phi_i^{\,n}\big)=\sum_{f\in\partial\Omega_i}A_f\,J_f+V_i\,S_i
  \label{eq:5-18}
\end{equation}
其中 $V_i$ 为第 $i$ 个控制体的体积、$A_f$ 为界面面积、$J_f$ 为经该界面的数值通量、$S_i$ 为该控制体内的源项；同一公共界面在两个相邻控制体上取同一个 $J_f$，故内界面通量在求和时两两抵消、只剩最外层边界项。这就是“守恒性在代数层面严格成立且与网格是否均匀无关”的代数含义：它不依赖网格疏密与是否均匀，只依赖“界面通量单值”这一条。

求解代价与求解效率可先作量级判断。元体平衡导出的半离散系统在径向上只有相邻节点耦合，故每步需求解一个三对角线代数系统（对角占优性见 5.1.3），追赶法使每步代价为 $O(N)$、与网格点数成正比且不随带宽增长；整个问题的总代价因此为式~\eqref{eq:5-19}

\begin{equation}
  W_{\rm tot}=O(N\cdot N_t),\qquad N_t=\frac{t_{\rm end}}{\Delta t_{\rm sub}}
  \label{eq:5-19}
\end{equation}
其中 $N$ 为径向区间数、$N_t$ 为内部子步总数。四个小问的实测墙钟时间（单机单线程，均取自数值口径总表的登记值）为：问题一单遍约 45 s（含结果写出与断言约 55 s）；问题二 9.9 s——求解核经数学等价的实现级优化后由 47.6 s 降至 9.9 s（加速 4.80 倍），关键量逐位一致；问题三 248.2 s；问题四 487 s，约 8.1 min。问题四最慢，与其径向区间数加倍（160）且时程最长、内部迭代次数最多相符。本文认为该量级已完全满足工程可用性，故不再以并行化或改换编译型语言等方式压缩，以保持交付实现单一、可复现。时间上问题一温度场采用时间方向精确推进，其余场量全部采用后向欧拉全隐式并配合追赶法，理由是隐式格式无条件稳定、长时程计算不受步长限制，而问题一温度场在半离散后恰为常系数线性系统、积分可闭式求得，可直接消除时间方向的离散误差。检验手段详见第六章；本文各环节的选型见表~\ref{tab:5-1}。

进一步说明排除口径。空间上不采用二阶中心差分：其中心点须作特殊处理，且守恒性只在代数上近似成立，而元体平衡在逐控制体意义下严格守恒、与网格是否均匀无关。时间上不采用显式格式：扩散算子的稳定性限幅会把含水率侧的子步压至 $10^{-3}$ s 量级，长时程推进的计算量不可接受，而全隐式无条件稳定、长时程计算不受步长限制。温度场亦不采用“显式格式加大子步”的逼近路线：该路线的时间离散误差仅随步长线性下降，而问题一温度场在半离散后恰为常系数线性系统、积分可闭式求得，可直接把时间离散误差归零（步长 1 s 与 60 s 的结果相差 $3.1\times10^{-11}$ ℃，即为其量化对照）。



\settabw{6}
\begin{longtable}{>{\raggedright\arraybackslash}m{\dimexpr 0.174\tblw\relax}>{\raggedright\arraybackslash}m{\dimexpr 0.306\tblw\relax}>{\raggedright\arraybackslash}m{\dimexpr 0.520\tblw\relax}}
  \caption{求解方案选型对比}\label{tab:5-1}\\
  \toprule
    环节 & 本文选用（备选） & 选用依据与量化对照 \\
  \midrule
  \endfirsthead
  \multicolumn{3}{r}{\footnotesize（续）}\\
  \toprule
    环节 & 本文选用（备选） & 选用依据与量化对照 \\
  \midrule
  \endhead
  \bottomrule
  \endfoot
  \bottomrule
  \endlastfoot
  空间离散 & 元体平衡（单元中心有限体积） & 逐控制体严格守恒；中心半格使中心奇点无须弱形式；中心系数两路推导差 $2.8\times10^{-14}$ K \\
  温度场时间推进 & 精确推进（显式／后向欧拉） & 常系数线性系统积分可闭式；步长 1 s 与 60 s 相差 $3.1\times10^{-11}$ ℃，时间离散误差归零 \\
  含水率场时间推进 & 后向欧拉全隐式（显式／梯形） & 三对角严格对角占优，无条件稳定且满足离散极值原理；子步 0.1 s 下误差 $10^{-5}$ kg/kg \\
  子步与耦合 & 交替推进＋固定子步（顺序单向／自适应） & 物性双向依赖，单向求解误差沿时间累积；自适应轨与固定子步数值等价（不超过 $9.3\times10^{-6}$ kg/kg） \\
  界面变系数 & 沿含水率的积分平均（算术／调和平均） & 由界面通量守恒推出；尾部胞内 $D$ 跨约两个数量级，调和平均使时长高估 $52\%$ 以上，仅作对照臂 \\
  移动边界 & 物料坐标固定域化（固定坐标＋插值） & 表面恒为节点、输出点恒为计算点，消除移动域插值误差与表面列收敛困难，并揭示压缩温升 \\
\end{longtable}

\subsection{问题一：预热平衡阶段（常物性 · 完全解耦）}
\subsubsection{分析与建模思路}
药材在 $0$–$1800$ s 内处于预热平衡阶段：环境温度逐步上升，水分尚未大量脱除。首要判断是是否需要分布式求解，判据是两个毕奥数：

\begin{equation}
  Bi_T=\frac{hR_0}{k}\approx1.39,\qquad Bi_m=\frac{k_mR_0}{D}\approx3.24
  \label{eq:5-14}
\end{equation}
二者均为一阶量级，表明内部导热阻与外部对流阻同量级，药材内部必然存在显著的径向梯度，集总参数法不适用\textsuperscript{\cite{ref1,ref8}}。

其次是能否用半无限介质近似：扩散域厚度与半径之比不可忽略，忽略曲率会低估径向扩散路径。自研对照显示，同一条件下半无限近似把表面含水率高估约 $6.7\%$（$1.6113$ 对 $1.5104$ kg/kg），故必须保留曲率项\textsuperscript{\cite{ref1}}。

最后是两场能否解耦：问题一物性按附录 2 取常值，扩散系数只依赖含水率、不含温度，含水率方程也不含温度，两场无反向依赖，故可严格解耦依次求解——这也是后文对温度场使用时间方向精确推进的前提。

\subsubsection{b 数据说明}
环境温度与水分浓度取自附件 1，共 $241$ 个实测点（$0$–$14400$ s，间隔 $60$ s），缺失点 $0$ 个、三倍标准差离群点 $0$ 个；节点间采用线性插值、不做平滑（对照计算表明平滑与否对结果的影响不超过 $1.08\times10^{-2}$ ℃）。初始温度 $28$ ℃、初始干基含水率 $2.55$ kg/kg 由题面给定。问题一的整个计算窗口 $0$–$1800$ s 完全落在升温段内，边界数据全部取自实测，不涉及任何外推。

\subsubsection{控制方程与定解条件}
问题一的物性取附录 2 常物性，控制方程为

\begin{equation}
  \rho c_p\frac{\partial T}{\partial t}=\frac{k}{r}\frac{\partial}{\partial r}\Big(r\frac{\partial T}{\partial r}\Big),\qquad \frac{\partial C}{\partial t}=\frac{1}{r}\frac{\partial}{\partial r}\Big(D(C)\,r\frac{\partial C}{\partial r}\Big)
  \label{eq:5-2}
\end{equation}
其中扩散系数仍随含水率变化，这一依赖关系是问题一唯一的非线性来源，但它只出现在含水率方程内部，不与温度方程交叉。定解条件为式~\eqref{eq:5-3} 至式~\eqref{eq:5-5}

\begin{equation}
  T(r,0)=T_0=28\ \mathrm{^\circ C},\qquad C(r,0)=C_0=2.55\ \mathrm{kg/kg}
  \label{eq:5-3}
\end{equation}
\begin{equation}
  \frac{\partial T}{\partial r}\Big|_{r=0}=0,\qquad \frac{\partial C}{\partial r}\Big|_{r=0}=0
  \label{eq:5-4}
\end{equation}
\begin{equation}
  -k\frac{\partial T}{\partial r}\Big|_{r=R_0}=h\big[T(R_0,t)-T_\infty(t)\big],\qquad -D\frac{\partial C}{\partial r}\Big|_{r=R_0}=k_m\big[C(R_0,t)-C_\infty(t)\big]
  \label{eq:5-5}
\end{equation}
其中 $h=25$ W/(m\textsuperscript{2}·K)、$k_m=8\times10^{-7}$ m/s 由附录 2 给出并全程沿用；$T_\infty(t)$、$C_\infty(t)$ 由附件 1 线性插值给出。中心条件 (5-4) 来自轴对称性，它使中心成为自然边界而非物理边界；表面条件 (5-5) 为第三类（对流换热与对流传质）条件，即表面通量由流体侧对流决定。两个条件的方向可作自检：预热阶段表面升温，故表面温度梯度为正向，与式~\eqref{eq:5-5} 右端符号一致。

\subsubsection{离散格式与求解算法}
空间离散采用元体平衡（单元中心有限体积），取径向区间数 $N=80$（$\Delta r$ 取 0.25 mm、$81$ 个节点），理由是使题面要求的 $0.1$ cm 输出点恰为计算点，从源头免除输出插值误差。中心控制体取半格，中心奇点由两条相互独立的路径装配（洛必达极限／半格通量平衡），系数完全一致、解级比对差异为 $2.8\times10^{-14}$ K；本文以半格通量平衡为交付路径，另一条留作交叉验证。该半格平衡可显式写出：第 $0$ 号控制体自中心 $r=0$ 延伸至第一个界面 $r=\Delta r/2$，体积为 $\pi(\Delta r/2)^2$，而中心侧面积为 $0$，故其平衡式只留外侧一项，中心奇点不出现在式中：

\begin{equation}
  \frac{\pi(\Delta r/2)^2}{\Delta t}\big(\phi_0^{\,n+1}-\phi_0^{\,n}\big)=2\pi\frac{\Delta r}{2}\,J_{1/2}
  \label{eq:5-20}
\end{equation}
其中 $J_{1/2}$ 为中心与第一个节点之间的界面数值通量。该式说明中心点不必作任何特殊近似：半格控制体把奇点化为一个“少一侧通量”的普通平衡式，这正是两条路径系数完全一致的代数原因。

温度场时间推进采用时间方向精确推进。常物性使半离散后的温度方程成为常系数线性系统：

\begin{equation}
  \frac{\mathrm d\mathbf T}{\mathrm dt}=\mathbf A\mathbf T+\mathbf w_0(t)+\mathbf w_1(t)
  \label{eq:5-6}
\end{equation}
其中 $\mathbf A$ 为半离散系数矩阵，为严格对角占优的非奇异 M 矩阵，故可逆；边界项由 $\mathbf w_0$、$\mathbf w_1$ 承载。附件 1 为等步长 $60$ s 数据且节点间线性插值，故边界在每个 $1$ s 子区间内严格分段线性，积分可闭式求得，得递推式~\eqref{eq:5-7}：

\begin{equation}
  \mathbf T^{n+1}=E\,\mathbf T^{n}+G_0\mathbf w_0^{n}+G_1\mathbf w_1^{n},\quad E=e^{\mathbf A h},\quad G_0=\mathbf A^{-1}(E-\mathbf I),\quad G_1=\mathbf A^{-2}(E-\mathbf I-\mathbf A h)
  \label{eq:5-7}
\end{equation}
$E$、$G_0$、$G_1$ 一经离线预分解即可全程复用。由于积分是闭式，该格式不引入任何时间离散误差；实测步长无关性为：分段步长 $h$ 取 $1$ s 与 $60$ s 的结果相差 $3.1\times10^{-11}$ ℃。

其适用边界为常系数前提：物性一旦随温度或含水率变化，闭式积分不再成立，故该推进只用于问题一，问题二至问题四的温度场改用后向欧拉全隐式。

含水率场时间推进采用后向欧拉全隐式。因扩散系数依赖含水率，方程为非线性，采用后向欧拉配合欠松弛迭代：

\begin{equation}
  \frac{\mathbf C^{n+1}-\mathbf C^{n}}{\Delta t}=\mathbf L_C\big(\mathbf C^{(k)}\big)\mathbf C^{(k+1)}+\mathbf g^{n+1},\qquad \omega=0.7
  \label{eq:5-8}
\end{equation}
其中上标 $(k)$ 为迭代层而非时间层；每步的新解与旧解按松弛因子加权组合：

\begin{equation}
  \mathbf C^{(k+1)}=(1-\omega)\,\mathbf C^{(k)}+\omega\,\tilde{\mathbf C}^{(k+1)},\qquad \omega=0.7
  \label{eq:5-15}
\end{equation}
迭代矩阵为三对角；由对角元大于同行两个非对角元之和可知其严格对角占优，而严格对角占优正是追赶法可用与极值原理成立的代数前提，故追赶法无条件稳定并满足离散极值原理，含水率在计算中不会越出初值与边界值所夹的范围。含水率子步取 $0.1$ s。精度预算由输出要求反推：四位小数输出下，空间离散误差不超过 $1.7\times10^{-5}$ ℃、时间误差量级仅 $10^{-5}$ kg/kg，均在最低有效位之下。

问题一的求解原理可概括为三段式：空间上以元体平衡离散、温度场沿时间方向精确推进、含水率场以全隐式配合欠松弛迭代。其求解流程如图~\ref{fig:5-3} 所示。


% 图5-3 ＝ F-07（F 号）／图49（PDF 文件索引号）／问题一求解原理示意图
\begin{figure}[htbp]
  \centering
\includegraphics[width=0.990\textwidth]{../../05\_成品图/图49-F07\_Q1求解原理示意图.pdf}
  \caption{问题一求解原理示意图}\label{fig:5-3}
\end{figure}

由图~\ref{fig:5-3} 可见，三段之间的接口是“先空间、再按场分推”，温度场与含水率场只在物性求值处发生单向联系；这意味着问题一的解耦不是额外的近似假设，而是“物性只依赖含水率”这一事实的直接结果，故可依次求解而不引入附加误差。

\subsubsection{结果与分析}
按题面格式给出 $0$–$1800$ s 内若干时刻、若干位置的结果，见表~\ref{tab:1} 与表~\ref{tab:2}。



\settabw{12}
\begin{longtable}{>{\centering\arraybackslash}m{\dimexpr 0.154\tblw\relax}>{\centering\arraybackslash}m{\dimexpr 0.169\tblw\relax}>{\centering\arraybackslash}m{\dimexpr 0.169\tblw\relax}>{\centering\arraybackslash}m{\dimexpr 0.169\tblw\relax}>{\centering\arraybackslash}m{\dimexpr 0.169\tblw\relax}>{\centering\arraybackslash}m{\dimexpr 0.170\tblw\relax}}
  \caption{30 分钟内药材的温度（单位：℃；各列为到药材中心的距离，单位：cm）}\label{tab:1}\\
  \toprule
    时间/s & 0 & 0.5 & 1.0 & 1.5 & 2.0 \\
  \midrule
  \endfirsthead
  \multicolumn{6}{r}{\footnotesize（续）}\\
  \toprule
    时间/s & 0 & 0.5 & 1.0 & 1.5 & 2.0 \\
  \midrule
  \endhead
  \bottomrule
  \endfoot
  \bottomrule
  \endlastfoot
  100 & 28.0001 & 28.0003 & 28.0041 & 28.0327 & 28.1800 \\
  300 & 28.0408 & 28.0635 & 28.1514 & 28.3680 & 28.8487 \\
  600 & 28.4534 & 28.5360 & 28.8039 & 29.3159 & 30.1652 \\
  900 & 29.3243 & 29.4583 & 29.8754 & 30.6161 & 31.7304 \\
  1200 & 30.5427 & 30.7098 & 31.2223 & 32.1125 & 33.4276 \\
  1500 & 31.9957 & 32.1867 & 32.7660 & 33.7463 & 35.1203 \\
  1800 & 33.5753 & 33.7720 & 34.3642 & 35.3621 & 36.7855 \\
\end{longtable}

同一组输出时刻与同一组半径位置上的水分浓度见表~\ref{tab:2}，两表共用同一套输出契约与有效数字口径。



\settabw{12}
\begin{longtable}{>{\centering\arraybackslash}m{\dimexpr 0.167\tblw\relax}>{\centering\arraybackslash}m{\dimexpr 0.167\tblw\relax}>{\centering\arraybackslash}m{\dimexpr 0.167\tblw\relax}>{\centering\arraybackslash}m{\dimexpr 0.167\tblw\relax}>{\centering\arraybackslash}m{\dimexpr 0.167\tblw\relax}>{\centering\arraybackslash}m{\dimexpr 0.165\tblw\relax}}
  \caption{30 分钟内药材的水分浓度（单位：kg/kg；各列为到药材中心的距离，单位：cm）}\label{tab:2}\\
  \toprule
    时间/s & 0 & 0.5 & 1.0 & 1.5 & 2.0 \\
  \midrule
  \endfirsthead
  \multicolumn{6}{r}{\footnotesize（续）}\\
  \toprule
    时间/s & 0 & 0.5 & 1.0 & 1.5 & 2.0 \\
  \midrule
  \endhead
  \bottomrule
  \endfoot
  \bottomrule
  \endlastfoot
  100 & 2.5500 & 2.5500 & 2.5500 & 2.5500 & 2.2490 \\
  300 & 2.5500 & 2.5500 & 2.5500 & 2.5492 & 2.0517 \\
  600 & 2.5500 & 2.5500 & 2.5500 & 2.5352 & 1.8774 \\
  900 & 2.5500 & 2.5500 & 2.5497 & 2.5045 & 1.7550 \\
  1200 & 2.5500 & 2.5500 & 2.5482 & 2.4646 & 1.6588 \\
  1500 & 2.5500 & 2.5499 & 2.5445 & 2.4206 & 1.5789 \\
  1800 & 2.5500 & 2.5497 & 2.5383 & 2.3755 & 1.5104 \\
\end{longtable}

温度场与水分浓度场的时空演化如图~\ref{fig:5-5} 与图~\ref{fig:5-6} 所示。

该图横轴为到药材中心的距离、纵轴为时间、色标为温度；由图~\ref{fig:5-5} 可见，径向上始终维持由表及里的温度梯度，热量在 1800 s 内尚未渗透至中心。


\begin{figure}[htbp]
  \centering
% 图5-5 ＝ F-08（F 号）／图06（PDF 文件索引号）／问题一温度场时空分布
  \begin{minipage}[t]{0.485\textwidth}
    \centering
    \includegraphics[width=1.000\linewidth]{../../05\_成品图/图06-Q1温度场时空分布.pdf}
    \caption{问题一温度场时空分布}\label{fig:5-5}
  \end{minipage}
  \hfill
% 图5-6 ＝ F-09（F 号）／图07（PDF 文件索引号）／问题一水分浓度场时空分布
  \begin{minipage}[t]{0.485\textwidth}
    \centering
    \includegraphics[width=1.000\linewidth]{../../05\_成品图/图07-Q1水分浓度场时空分布.pdf}
    \caption{问题一水分浓度场时空分布}\label{fig:5-6}
  \end{minipage}
\end{figure}

该图与图~\ref{fig:5-5} 同构，色标为干基含水率；由图~\ref{fig:5-6} 可见，含水率的变化带厚度约为 0.5–1.0 cm，其作用范围远小于同期的温度场。

由表~\ref{tab:1}、表~\ref{tab:2} 与图~\ref{fig:5-5}、图~\ref{fig:5-6}，可归纳出三条规律。

其一，温度外高内低，且 $1800$ s 内热量尚未渗透至中心。1800 s 时表面温度 $36.7855$ ℃、中心温度 $33.5753$ ℃，由表~\ref{tab:1} 各行沿径向单调可见；$1800$ s 约为热扩散特征时间 $\tau_T\approx2.37\times10^{3}$ s 的 $0.76$ 倍，故预热阶段远未结束。

其二，水分变化仅局限在表层。表~\ref{tab:2} 显示中心含水率在四位小数下始终保持 $2.5500$，即表层微降、内部尚未启动；而表面含水率由 $2.55$ 降至 $1.5104$ kg/kg，降幅 $40.8\%$，变化带厚度约为 $0.5$–$1.0$ cm。这意味着水分脱除过程尚处在最外层。

其三，两个过程的时间尺度相差约 $34$ 倍，且温差仍在建立过程中。两场特征时间由各自扩散尺度定义：

\begin{equation}
  \tau_T=\frac{R_0^2}{\alpha}=\frac{\rho c_pR_0^2}{k}\approx2.37\times10^{3}\ \mathrm s,\qquad \tau_C=\frac{R_0^2}{D}\approx8.1\times10^{4}\ \mathrm s
  \label{eq:5-16}
\end{equation}
二者之比约 $34$，而表面与中心的温差随时间持续增大。这意味着“温度平衡快、水分脱除慢”是本题的固有结构，直接决定后续长时程干燥的耗时量级，而预热段的温度梯度尚未进入准稳态。

上述规律的共同根源是热扩散与质扩散的时间尺度相差约 34 倍（即 $\tau_T$ 与 $\tau_C$ 之比），而水分的作用层厚度由扩散长度决定而非由药材半径决定：

\begin{equation}
  \ell_D(t)=\sqrt{Dt}
  \label{eq:5-17}
\end{equation}
故该时段内部未变化是物理必然，不是数值误差或模型缺陷；且在第三类对流边界下径向温差很快饱和、整体温度随环境平缓上移，故问题一全程处于预热而非平衡阶段。

问题一的参数灵敏度结构由第六章 6.6 与 6.7 集中给出（局部单因素与全局两族方法各一节），正文不再单独呈现该图件。局部灵敏度结果显示，环境温度与扩散系数为主控因素，全部算例中结果随参数单调变化、未出现临界敏感或反号现象；图~\ref{fig:5-5} 与图~\ref{fig:5-6} 的时空分布也从侧面支持这一判断。

本问的中间结果以网格加密序列的形式留档，见表~\ref{tab:5-3}。本文认为，把加密序列与对拍偏差并列给出，比只报一个末值更能说明“报告值已收敛”这一判断的依据；图~\ref{fig:5-5} 与图~\ref{fig:5-6} 的时空分布也从另一侧面印证了同一结论。



\settabw{8}
\begin{longtable}{>{\centering\arraybackslash}m{\dimexpr 0.500\tblw\relax}>{\centering\arraybackslash}m{\dimexpr 0.146\tblw\relax}>{\centering\arraybackslash}m{\dimexpr 0.146\tblw\relax}>{\centering\arraybackslash}m{\dimexpr 0.208\tblw\relax}}
  \caption{问题一：网格加密序列与解析对拍对照}\label{tab:5-3}\\
  \toprule
    径向网格 Δr/mm & 1.0 & 0.5 & 0.25 \\
  \midrule
  \endfirsthead
  \multicolumn{4}{r}{\footnotesize（续）}\\
  \toprule
    径向网格 Δr/mm & 1.0 & 0.5 & 0.25 \\
  \midrule
  \endhead
  \bottomrule
  \endfoot
  \bottomrule
  \endlastfoot
  与圆柱级数解的对拍偏差 / 10\textsuperscript{-4} ℃ & 3.886 & 0.975 & 0.244 \\
  含水率侧对拍偏差 / 10\textsuperscript{-3} kg/kg & 1.029 & 0.267 & 0.083 \\
  温度场观测阶（中心／表面／均方根） & — & — & 2.002／2.006／2.003 \\
  中心网格收敛指数 & — & — & ≤0.0001\% \\
\end{longtable}

注：对拍偏差随网格二阶下降，说明残差已归于空间离散；时间侧分段步长取 1 s 与 60 s 的结果相差 $3.1\times10^{-11}$ ℃，即时间离散误差已归零。

\subsubsection{模型检验（要点）}
问题一的检验分四组：网格与时间收敛性（空间二阶，中心网格收敛指数不超过 $0.0001\%$）、圆柱级数解析解对拍、第二套独立实现互验、物理自检与守恒核算；另有数据预处理对照与边界数据不确定度传播两项，见第六章 6.8 与 6.11。三档网格与三档步长的收敛序列见图~\ref{fig:A-1}，与圆柱级数解的对拍结果见图~\ref{fig:A-2}；温度场与含水率场的时空分布可回看 图~\ref{fig:5-5} 与 图~\ref{fig:5-6}，用以对照本节所报的径向量值。证据与图表统一见第六章，本节不重复其数值。

\subsection{问题二：整个烘干过程（变物性 · 双向强耦合）}
\subsubsection{分析与建模思路}
与问题一相比，变化集中在两处：物性体系由附录 2 换为附录 3，此时密度、比热容、导热系数与扩散系数均为含水率与温度的函数；耦合方式由严格解耦转为双向强耦合。

附录 3 的四项物性均为状态量的经验函数，扩散系数为两个独立负指数因子之积（含温度因子）；因该因子随温度升高显著提高扩散系数，这正是温度到水分这一耦合通道的物理来源。三套物性在含水率 2.55 kg/kg、温度 301.15 K 处的取值对照为：导热系数 0.36／0.4830／0.2637 W/(m·K)，扩散系数 4.9377／5.6417／1.0471×10\textsuperscript{-9} m\textsuperscript{2}/s，密度 820／976.4／989.5 kg/m\textsuperscript{3}，比热容 2600／3415.3／3394.4 J/(kg·K)（依附录 2／3／4 顺序）；三者随含水率与温度的变化规律及其影响，已并入第六章 6.6 与 6.7 的灵敏度与耦合强度分析，正文不再单独呈现演化曲线图。

耦合的机理是：水分扩散系数随温度升高而增大，而热物性随含水率下降而改变。若先算温度再算含水率而不回代，物性将始终取旧状态，误差会沿时间累积并以系统性偏差进入结果。本文采用交替推进并令物性在每个内部子步按局部状态更新，以化解这一矛盾\textsuperscript{\cite{ref9,ref10}}。四项物性中，密度与比热容随含水率接近线性，而导热系数与扩散系数呈强非线性（其初始状态取值见本节文字），这正是后续耦合强度量化与灵敏度分析的焦点，定量结果见第六章 6.6 与 6.7；图~\ref{fig:5-10} 给出的双向耦合回路即由此而来。

两场双向耦合的具体回路如图~\ref{fig:5-10} 所示。


% 图5-10 ＝ F-21（F 号）／图50（PDF 文件索引号）／问题二双向耦合关系图
\begin{figure}[htbp]
  \centering
\includegraphics[width=0.800\textwidth]{../../05\_成品图/图50-F21\_Q2双向耦合关系图.pdf}
  \caption{问题二双向耦合关系图}\label{fig:5-10}
\end{figure}

由图~\ref{fig:5-10} 可见，温度经扩散系数影响含水率、含水率又经热物性反向影响温度，两条通路共同构成本问的双向依赖；这意味着“先温度后含水率”的单向顺序求解会在每个子步留下未被回代的状态偏差，故必须按子步更新物性并反复推进。

\subsubsection{定解条件的变化}
初始条件与问题一相同。环境数据在 14400 s 之后不再有实测覆盖，本文按附件 1 末段实测均值外推为恒定值：环境温度 50.00 ℃、环境水分浓度 0.04999 kg/kg；该外推自问题二起生效。对恒温段实测点的线性趋势检验显示斜率不显著、与零无显著差异，故以段均值作常值延拓不引入趋势性偏差；其影响在第六章以灵敏度与扰动方式单独评估。

\subsubsection{离散格式与求解算法}
方程形式与问题一保持一致（仍取散度形式~\eqref{eq:5-1}），变化在物性求值方式与耦合推进方式。

界面变系数的取法是本节的关键。控制体交界面上导热系数与扩散系数不再连续，本文取沿含水率的积分平均（式~\eqref{eq:5-9}）：

\begin{equation}
  \langle D\rangle=\frac{1}{\Delta C}\int_{C_1}^{C_2}D\,\mathrm dC
  \label{eq:5-9}
\end{equation}

该式由界面通量守恒推出而非格式惯例。干燥尾部单个控制体内扩散系数跨越约两个数量级，取调和平均会严重低估有效扩散（问题三侧量化：取法错误可使烘干时长被高估 52\% 以上）；故本文以积分平均为准，调和平均仅作对照臂。

该取法的必要性有其边界：界面两侧系数相差不大时两法在四位小数下不可辨识（问题一侧最大差 $8.5\times10^{-6}$ kg/kg），故其收益集中于干燥尾部的高湿差段。

% 图5-11 ＝ F-22（F 号）／图19（PDF 文件索引号）／交替推进流程图
%   ★ 本轮（用户第 8 项）：本块**下移两段**——原位置使 wrapfigure 的预留区落在 p13、图却画到
%     p14，导致 p13 出现"为图留白但图不在"的空档。下移后预留区与图同在 p14。
\begin{wrapfigure}{r}{0.519\textwidth}
  \centering
  \includegraphics[width=\linewidth]{../../05\_成品图/图19-IMEX交替推进算法流程图.pdf}
  \vspace{2pt}
  \captionof{figure}{交替推进流程图}\label{fig:5-11}
\end{wrapfigure}

时间推进采用交替推进：每个内部子步先解温度（全隐式，物性冻结于子步起点），再解含水率（欠松弛迭代，$\omega$ 取 0.7、容限取 10\textsuperscript{-10}），物性在每个内部子步按局部状态更新，迭代收敛性的理论依据见文献\textsuperscript{\cite{ref11}}。物性冻结造成的滞后正比于子步长度，子步取 1/32 s 使该滞后不进入四位小数量级。“按局部状态更新”可写成显式形式：设第 $n$ 个子步起点状态为 $(C^n,T^n)$，则该子步内取

\begin{equation}
  \begin{split}
  k^{n} &= k\big(C^{n},T^{n}\big),\\
  D^{n} &= D\big(C^{n},T^{n}\big),\\
  (\rho c_p)^{n} &= (\rho c_p)\big(C^{n},T^{n}\big)
  \end{split}
  \label{eq:5-21}
\end{equation}
即物性在子步起点冻结、于下一步刷新。由该式可见，冻结引入的滞后量与子步长度成正比（量级为 $\partial_t k\cdot\Delta t_{\rm sub}$），这就是子步须取到 1/32 s 的直接原因。

上述交替推进的完整流程如图~\ref{fig:5-11} 所示。

由图~\ref{fig:5-11} 可见，每个内部子步按“更新物性、解温度、解含水率、判收敛”的顺序推进，物性在子步起点冻结、于下一步刷新；这意味着格式的显隐分离（两场全隐、物性显式求值）使误差来源可逐项归因，而不混入未离散的状态量。

\subsubsection{结果与分析}
按题面格式给出 0–3 h 内每隔 0.5 h、若干位置的结果，见表~\ref{tab:3} 与表~\ref{tab:4}。



\settabw{12}
\begin{longtable}{>{\centering\arraybackslash}m{\dimexpr 0.154\tblw\relax}>{\centering\arraybackslash}m{\dimexpr 0.169\tblw\relax}>{\centering\arraybackslash}m{\dimexpr 0.169\tblw\relax}>{\centering\arraybackslash}m{\dimexpr 0.169\tblw\relax}>{\centering\arraybackslash}m{\dimexpr 0.169\tblw\relax}>{\centering\arraybackslash}m{\dimexpr 0.170\tblw\relax}}
  \caption{3 小时内药材的温度（单位：℃；各列为到药材中心的距离，单位：cm）}\label{tab:3}\\
  \toprule
    时间/h & 0 & 0.5 & 1.0 & 1.5 & 2.0 \\
  \midrule
  \endfirsthead
  \multicolumn{6}{r}{\footnotesize（续）}\\
  \toprule
    时间/h & 0 & 0.5 & 1.0 & 1.5 & 2.0 \\
  \midrule
  \endhead
  \bottomrule
  \endfoot
  \bottomrule
  \endlastfoot
  0.5 & 32.1893 & 32.3822 & 32.9660 & 33.9608 & 35.4131 \\
  1.0 & 40.3816 & 40.5536 & 41.0605 & 41.8783 & 42.9977 \\
  1.5 & 45.8468 & 45.9348 & 46.1930 & 46.6051 & 47.1401 \\
  2.0 & 48.4502 & 48.4881 & 48.5984 & 48.7736 & 49.0033 \\
  2.5 & 49.4670 & 49.4792 & 49.5136 & 49.5653 & 49.6609 \\
  3.0 & 49.8495 & 49.8553 & 49.8746 & 49.9101 & 49.9664 \\
\end{longtable}

与表~\ref{tab:3} 同组时刻、同组位置上的水分浓度见表~\ref{tab:4}，其行次与题面模板一致。



\settabw{12}
\begin{longtable}{>{\centering\arraybackslash}m{\dimexpr 0.167\tblw\relax}>{\centering\arraybackslash}m{\dimexpr 0.167\tblw\relax}>{\centering\arraybackslash}m{\dimexpr 0.167\tblw\relax}>{\centering\arraybackslash}m{\dimexpr 0.167\tblw\relax}>{\centering\arraybackslash}m{\dimexpr 0.167\tblw\relax}>{\centering\arraybackslash}m{\dimexpr 0.165\tblw\relax}}
  \caption{3 小时内药材的水分浓度（单位：kg/kg；各列为到药材中心的距离，单位：cm）}\label{tab:4}\\
  \toprule
    时间/h & 0 & 0.5 & 1.0 & 1.5 & 2.0 \\
  \midrule
  \endfirsthead
  \multicolumn{6}{r}{\footnotesize（续）}\\
  \toprule
    时间/h & 0 & 0.5 & 1.0 & 1.5 & 2.0 \\
  \midrule
  \endhead
  \bottomrule
  \endfoot
  \bottomrule
  \endlastfoot
  0.5 & 2.5499 & 2.5489 & 2.5255 & 2.3258 & 1.6486 \\
  1.0 & 2.5257 & 2.4947 & 2.3578 & 2.0230 & 1.4711 \\
  1.5 & 2.3860 & 2.3256 & 2.1344 & 1.8020 & 1.3475 \\
  2.0 & 2.1708 & 2.1084 & 1.9236 & 1.6259 & 1.2311 \\
  2.5 & 1.9567 & 1.9006 & 1.7360 & 1.4720 & 1.1166 \\
  3.0 & 1.7662 & 1.7165 & 1.5702 & 1.3333 & 1.0081 \\
\end{longtable}

注：表~\ref{tab:3}、表~\ref{tab:4} 的行次按题面模板，自 0.5 h 起，故不列 0 h 行。

由表~\ref{tab:3} 与表~\ref{tab:4}，可归纳出两条结论。

其一，温度已基本趋于环境值，而含水率仍远高于达标阈值。3 h 时中心温度 49.8495 ℃、表面温度 49.9664 ℃，均已逼近环境温度 50.00 ℃；而中心含水率 1.7662 kg/kg、表面含水率 1.0081 kg/kg，较初值的降幅分别为 30.7\% 与 60.5\%。这意味着干燥远未完成，与问题一揭示的两尺度分离相互印证。

其二，两层结果在重叠段的差异是模型必然，不是计算错误，而耦合强度可定量描述。变物性口径下的热扩散率较常物性低 14.2\%（以问题一为基准），导致同一时刻的温度与含水率出现系统性差异，这意味着问题一与问题二的结果本就不应完全一致；同时两场之间的相互影响在干燥中段达到最强（耦合强度已按“温度→水分”与“物性演化”两条通道量化），这意味着交替迭代在该阶段的价值最大。

其后的机理需更正一处常见误解：第三类传质边界下表面通量由传质系数与浓度差决定、与扩散系数无关，扩散系数只决定内部剖面的平坦程度；故同一时刻问题二表面含水率高于问题一，原因是其剖面更平坦而非传质更快。

上述差异的来源可作分解如图~\ref{fig:5-12} 所示。


% 图5-12 ＝ F-19（F 号）／图17（PDF 文件索引号）／问题一与问题二重叠段差异归因
\begin{wrapfigure}{r}{0.480\textwidth}
  \centering
  \includegraphics[width=\linewidth]{../../05\_成品图/图17-Q1Q2重叠段差异归因.pdf}
  \vspace{2pt}
  \captionof{figure}{问题一与问题二重叠段差异归因}\label{fig:5-12}
\end{wrapfigure}

该图给出两问在 0–1800 s 重叠段内中心与表面量的差异及其成因分解；由图~\ref{fig:5-12} 可见，变物性口径下的热扩散率较常物性低 14.2\%（以问题一为基准），与本节所报差异的量级一致。

本问耦合推进的设置与收敛性对照见表~\ref{tab:5-4}；图~\ref{fig:5-10} 给出的双向耦合回路正是该表“逐子步更新物性”一行的直接依据。



\settabw{6}
\begin{longtable}{>{\raggedright\arraybackslash}m{\dimexpr 0.237\tblw\relax}>{\raggedright\arraybackslash}m{\dimexpr 0.337\tblw\relax}>{\raggedright\arraybackslash}m{\dimexpr 0.426\tblw\relax}}
  \caption{问题二：耦合推进设置与收敛性对照}\label{tab:5-4}\\
  \toprule
    项 & 交付设置 & 对照或依据 \\
  \midrule
  \endfirsthead
  \multicolumn{3}{r}{\footnotesize（续）}\\
  \toprule
    项 & 交付设置 & 对照或依据 \\
  \midrule
  \endhead
  \bottomrule
  \endfoot
  \bottomrule
  \endlastfoot
  时间推进 & 交替推进（两场全隐、物性显式求值） & 温度经扩散系数影响含水率，含水率经热物性反作用于温度 \\
  内部子步 & 1/32 s & 物性冻结滞后与子步长度成正比 \\
  欠松弛因子 & 0.7 & 非线性内层迭代的收敛加速 \\
  内层迭代容限 & 10\textsuperscript{-10} & 保证每个子步真正收敛 \\
  物性更新时机 & 每个内部子步按局部状态 & 单向顺序求解会在每步留下未回代的状态偏差 \\
  自适应子步（对照臂） & 子步数降至固定子步的 18.75\% & 关键量偏差不超过 $9.3\times10^{-6}$ kg/kg \\
\end{longtable}

注：交付值取固定子步口径以保证可复现；自适应方案仅作对照臂，其收益与代价见表内对应行。

\subsubsection{模型检验（要点）}
问题二的检验以退化对照为核心（含强制解耦退化、独立实现互验、守恒核算、潜热对照与自适应步长对照）：退化对拍只证明框架与实现在常物性极限下正确，不等同于变物性物理本身正确，后者由灵敏度分析与参数扰动共同支撑。证据见第六章。

\subsection{问题三：烘干所需时间（反演 · 事件驱动终止）}
\subsubsection{分析与建模思路}
与问题二相比，本问仅动两处：时程由固定 3 h 改为事件驱动（达到判据即停），并新增判据要求全域干基含水率严格小于 0.15 kg/kg。

第一个关键判断是判据可否只在中心监测。由于中心到表面的扩散路径最长，且表面通量最先影响外层，中心恒为最慢点：任一时刻含水率沿半径自表面向中心单调递增，故“全域达标”与“中心达标”等价，判据只需在中心监测即可。该等价性可写成式~\eqref{eq:5-22}

\begin{equation}
  \max_{r\in[0,R_0]}C(r,t)=C(0,t)\qquad\Longrightarrow\qquad \Big[\max_{r}C(r,t)<0.15\ \Longleftrightarrow\ C(0,t)<0.15\Big]
  \label{eq:5-22}
\end{equation}
左端的等号由径向单调性给出，右端是判据的等价改写。故本文只在中心布置判据，而这一做法不减少任何严格性：全域判据与中心判据在数学上同解。

第二个关键判断是几何：本题设定下半径全程固定（收缩属问题四），故模型与离散可完全复用问题二，本问实质上是在既有求解器上加装终止机制。这一“复用加终止”的定位也决定了检验重点：既要验证终止机制本身，也要验证它与问题二在重叠段的逐位一致性。

\subsubsection{定解条件的变化}
终端判据取干基含水率严格小于 0.15 kg/kg，即等号处不触发；物性仍用附录 3，与问题二一致。需声明的是，本问把含水率区间由 1.0 kg/kg 以上扩展到 0.15 kg/kg，扩散系数经验式在低含水率端的外推幅度显著加大，外推段对烘干时长的贡献达 43\%–66\%；该代价已在第六章量化，并列为第七章的结论适用边界。另需界定本问所用“反演”一词的含义：此处指在严格判据下定位达标时刻，属单调、良态问题；它与扩散系数、辐射系数一类的病态参数反演性质不同，后者的统计反演方法与反问题正则化处理可参见文献\textsuperscript{\cite{ref19,ref20}}。

\subsubsection{离散格式与求解算法}
方程、边界与物性与问题二完全一致，沿用式~\eqref{eq:5-1} 与式~\eqref{eq:5-9}。本问新增两项机制：每个输出步检测全域含水率最大值是否小于 0.15 kg/kg；一旦触发，则在末两个输出步之间做子步级二分，定位达标时刻，精度取 10\textsuperscript{-3} s，即

\begin{equation}
  \max_r C(\cdot,t_{\rm end})<0.15,\qquad |t_{\rm end}-t^{*}|<10^{-3}\ \mathrm s
  \label{eq:5-10}
\end{equation}
内部步长不因输出变稀而放大：判据检测发生在输出步上，但时间推进仍保持 1/32 s 内部子步，以免牺牲精度、也免使二分区间内的时间分辨率被输出间隔绑架。问题三的事件驱动求解流程如图~\ref{fig:5-15} 所示。

由图~\ref{fig:5-15} 可见，判据检测挂在输出步上、而二分定位发生在末两个输出步之间；这意味着达标时刻的定位精度由内部子步与二分容限共同决定（10\textsuperscript{-3} s），而不是由 6 h 的输出间隔决定，故表~\ref{tab:5} 的烘干结束行不因输出稀疏而变粗。

\subsubsection{结果与分析}
按题面格式给出每隔 6 h、若干位置的水分浓度，并给出烘干结束行，见表~\ref{tab:5}。



\settabw{12}
\begin{longtable}{>{\centering\arraybackslash}m{\dimexpr 0.203\tblw\relax}>{\centering\arraybackslash}m{\dimexpr 0.159\tblw\relax}>{\centering\arraybackslash}m{\dimexpr 0.159\tblw\relax}>{\centering\arraybackslash}m{\dimexpr 0.159\tblw\relax}>{\centering\arraybackslash}m{\dimexpr 0.159\tblw\relax}>{\centering\arraybackslash}m{\dimexpr 0.161\tblw\relax}}
  \caption{药材的水分浓度随时间和位置的变化（单位：kg/kg；各列为到药材中心的距离，单位：cm）}\label{tab:5}\\
  \toprule
    时间/h & 0 & 0.5 & 1.0 & 1.5 & 2.0 \\
  \midrule
  \endfirsthead
  \multicolumn{6}{r}{\footnotesize（续）}\\
  \toprule
    时间/h & 0 & 0.5 & 1.0 & 1.5 & 2.0 \\
  \midrule
  \endhead
  \bottomrule
  \endfoot
  \bottomrule
  \endlastfoot
  6 & 1.0170 & 0.9881 & 0.9015 & 0.7550 & 0.5328 \\
  12 & 0.4562 & 0.4432 & 0.4031 & 0.3298 & 0.1642 \\
  18 & 0.2989 & 0.2916 & 0.2687 & 0.2253 & 0.0845 \\
  24 & 0.2378 & 0.2328 & 0.2167 & 0.1856 & 0.0665 \\
  30 & 0.2059 & 0.2019 & 0.1892 & 0.1641 & 0.0598 \\
  36 & 0.1859 & 0.1826 & 0.1718 & 0.1504 & 0.0566 \\
  42 & 0.1721 & 0.1692 & 0.1597 & 0.1407 & 0.0548 \\
  48 & 0.1619 & 0.1592 & 0.1507 & 0.1335 & 0.0537 \\
  54 & 0.1539 & 0.1515 & 0.1437 & 0.1278 & 0.0530 \\
  \textbf{烘干结束时间} & \textbf{0.1500} & 0.1477 & 0.1402 & 0.1249 & \textbf{0.0526} \\
\end{longtable}

注：行次按题面表~\ref{tab:5} 模板，自 6 h 起，末行为烘干结束时间；该行对应达标时刻 207113.094 s，即 57.5314 h，其时间列为精确值，已由结果文件逐格取回。

满足全域干基含水率严格低于 0.15 kg/kg 的烘干所需时间为

\begin{equation*}
  t_{\rm end}=57.53\ \mathrm h\approx2.40\ \text{天}
\end{equation*}
经三档网格与三档时间步长核验，报告值带离散误差带 ±0.08 h；合成口径为空间网格收敛指数估计与时间保守上界取较大者\textsuperscript{\cite{ref21}}，两类不确定度分列报告、不相加。另以解析量级估计对该报告值作一次独立于数值实现的量级校验：下界为 25.8 h、上界为 97.7 h，本文报告值落在两界之内。长时程含水率演化如图~\ref{fig:5-16} 所示。


\begin{wrapfigure}{l}{0.450\textwidth}
  \centering
  \includegraphics[width=1.000\linewidth]{../../05\_成品图/图51-F51\_Q3求解流程图.pdf}
  \captionof{figure}{问题三求解流程图（事件驱动与子步级二分）}\label{fig:5-15}
\end{wrapfigure}

% ★ 本轮（图表更新后全量检查发现）：原先第二张图与题注同放在上面这个 wrapfigure 内，
%   实测**题注不渲染**（PDF 全文检索「问题三长时程含水率演化」为 0 命中）——
%   根因是 wrapfigure 一个环境承载两个题注不可靠（全篇仅此一处如此写法，其余环绕图均单题注）。
%   改为**分离式**：流程图留环绕，长时程演化图独立成常规浮动体，题注必定渲染。
\begin{figure}[htbp]
  \centering
  \includegraphics[width=0.620\textwidth]{../../05\_成品图/图29-Q3长时程含水率演化.pdf}
  \caption{问题三长时程含水率演化}\label{fig:5-16}
\end{figure}

给出中心与表面含水率在全程 57.53 h 内的演化曲线。由图~\ref{fig:5-16} 可见，中心曲线在干燥后期明显变缓，而表面曲线早早进入平台，两条曲线的间距随时间不断拉大。

由表~\ref{tab:5} 与图~\ref{fig:5-16}，可归纳出两条规律。

其一，干燥过程呈明显的两段式：前 24 h 中心含水率相对初值下降约九成，其后约 33.5 h 仅由 0.2378 降至 0.1500 kg/kg，即最后少量水耗去了大半时间，烘干时长由末段的慢过程决定。

其二，中心恒为最慢点，而表面含水率早早趋于平衡。由表~\ref{tab:5} 可见每一时刻含水率沿半径单调递增、越靠近中心越高，且表面在 24 h 后基本稳定；这意味着判据只需在中心监测，与 5.3.1 的分析一致，同时干燥后期的控制环节趋势上转向药材内部扩散，表面侧参数（换热系数、传质系数、环境浓度）的影响均很小。

机理上，扩散系数经验式中含 exp(−0.45/C) 形式的因子，随含水率下降而急剧减小，导致干燥在后期自我减速。这意味着烘干时长的瓶颈在低含水率段；工程上应以最慢点即中心排产，并优先控制温度与改善物料特性，而非单纯加大风量。

长时程还伴随一次主导机制转变：随含水率下降，传质由外部对流控制转为内部扩散控制，故达标时间不可由单一时间常数线性外推——这正是采用事件驱动终止配合子步级二分、以严格判据锁定达标时刻的原因。

本问的长时程收敛性与定位过程见表~\ref{tab:5-5}；表中所列定位精度与图~\ref{fig:5-15} 的事件驱动流程相互印证，值得指出的是，达标时刻的精度由内部子步与二分容限共同决定，而不受 6 h 输出间隔的影响。



\settabw{4}
\begin{longtable}{>{\raggedright\arraybackslash}m{\dimexpr 0.417\tblw\relax}>{\raggedright\arraybackslash}m{\dimexpr 0.583\tblw\relax}}
  \caption{问题三：长时程收敛性与达标时刻定位对照}\label{tab:5-5}\\
  \toprule
    口径 & 结果 \\
  \midrule
  \endfirsthead
  \multicolumn{2}{r}{\footnotesize（续）}\\
  \toprule
    口径 & 结果 \\
  \midrule
  \endhead
  \bottomrule
  \endfoot
  \bottomrule
  \endlastfoot
  空间网格 & 4 档加密，内部各点与表面列的相邻相对变化均不超过 $10^{-4}$ \\
  时间子步 & 1/32 s，经前后加密验证已收敛 \\
  空间观测阶／网格收敛指数 & 1.650 ／ 0.0011\% \\
  时间误差带 & 取保守上界 ±0.0802 h \\
  二分定位精度 & $10^{-3}$ s \\
  达标时刻（精确值） & 207113.094 s，即 57.5314 h \\
  与问题二重叠段（前 3 h） & 逐位一致 \\
  解析量级界（下界／上界） & 25.8 h ／ 97.7 h \\
\end{longtable}

注：本表各行均为检验与定位过程的中间结果，报告值与其误差带的合成口径见 6.1。

\subsubsection{模型检验（要点）}
问题三的检验以长时程重验与低含水率端外推评估为重点（后者对烘干时长的贡献达 43\%–66\%，为本问最大不确定性来源）；与问题二重叠段的逐位一致构成强交叉印证，其余检验项见第六章。

\subsection{问题四：计入尺寸收缩的烘干时长（移动边界）}
\subsubsection{分析与建模思路}
问题四在问题二、三的架构上引入两点实质变化：物性换用附录 4 体系\textsuperscript{\cite{ref5,ref6}}；药材半径随时间收缩（由附件 2 给出 $R(t)$，2.000→1.198 cm），扩散域随之移动。

必须做坐标变换的理由来自两类困难：固定坐标系下计算域随半径收缩会使网格节点与输出位置错位、输出值须在移动网格上插值；表面附近控制体被不断压缩、体积趋零，表面列收敛性恶化（只能达到一阶），而表面恰是通量最大处。

本文引入物料坐标，把收缩域映射为固定计算域：

\begin{equation}
  \xi=\frac{r}{R(t)},\qquad \xi\in[0,1]\qquad\Longleftrightarrow\qquad r=\xi R(t),\qquad \partial_r=\frac{1}{R(t)}\,\partial_\xi
  \label{eq:5-23}
\end{equation}
该变换的三个要素是：自变量由 $r$ 换为 $\xi$；物理半径与物料坐标之间为仿射关系；径向偏导算子按链式法则带上因子 $1/R(t)$，域移动的全部效应都进入该因子。

该变换带来两项收益：表面恒为节点，输出点恒为计算点\textsuperscript{\cite{ref15}}；且物料坐标下同一物质点始终对应同一节点，质量方程无需引入对流项即可描述域的移动。其适用边界与仿射收缩假设绑定：若收缩分配随局部含水率变化，节点对应不再是一一对应，须改用加权分配律——本文未做该对照，第七章不足 2 已列为改进方向。坐标变换的几何含义如图~\ref{fig:5-17} 所示。


% 图5-17 ＝ F-05（F 号）／图48（PDF 文件索引号）／物质坐标变换示意图
\begin{figure}[htbp]
  \centering
\includegraphics[width=0.728\textwidth]{../../05\_成品图/图48-F05\_物质坐标变换示意.pdf}
  \caption{物质坐标变换示意图}\label{fig:5-17}
\end{figure}

给出物理域与计算域的对应关系；由图~\ref{fig:5-17} 可见，物料坐标下的节点数不随时间变化，这意味着网格与输出点全程保持同构，是本问消除上述两类数值困难的关键。

\subsubsection{定解条件的变化}
收缩律 $R(t)$ 由附件 2 的节点线性插值给出（145 个实测点、间隔 1800 s），259200 s 之后取常值 1.198 cm。其插值式为

\begin{equation}
  R(t)=R_j+(t-t_j)\cdot\frac{R_{j+1}-R_j}{t_{j+1}-t_j},\qquad t\in[t_j,t_{j+1}],\qquad t_{j+1}-t_j=1800\ \mathrm s
  \label{eq:5-24}
\end{equation}
且 $R(0)=2.000$ cm、$R(t)\ge1.198$ cm；半径单调不增，这也是表~\ref{tab:6} 判定 1.5 cm 与 2.0 cm 两列越界的依据。本文采用仿射收缩假设（物质点半径按比例缩放）：附件 2 只给出半径的单一序列、药材内部各处的收缩分配未被规定，故按仿射处理，并在第三章假设中一并声明；含水率按干基定义，收缩不改变各处干物质的量，这一点在 5.4.3 的推导中自然体现。终端判据与问题三相同，仍为全域干基含水率严格小于 0.15 kg/kg。

本问输入、参数与输出均由题面或附录给定，不存在待优化决策变量，故不套用“目标函数加决策变量加约束”的优化模板，工作集中在离散格式与问题特化策略。

\subsubsection{方程在物料坐标下的形式与算法}
经坐标变换，质量方程在物料坐标下写为

\begin{equation}
  \frac{\partial \tilde C}{\partial t}=\frac{1}{\xi R^2}\frac{\partial}{\partial\xi}\Big(\xi D\,\partial_\xi \tilde C\Big)
  \label{eq:5-11}
\end{equation}
其中 $\tilde C(\xi,t)=C(r,t)$ 为物料坐标下的含水率。式~\eqref{eq:5-11} 为干基形式，不含对流项与几何压缩项；所以能写成这一形式，是因为变量本身以干基计量，脱水不改变干物质的质量基准。能量方程则必须取守恒形式，把几何压缩项自动包含在时间导数之内：

\begin{equation}
  \xi\,\frac{\partial}{\partial t}\Big[\rho c_p T R^2\Big]=\frac{\partial}{\partial\xi}\Big(k\,\xi\,\partial_\xi T\Big)
  \label{eq:5-12}
\end{equation}
上述处理属于移动边界问题的边界固定化方法族\textsuperscript{\cite{ref12,ref15,ref16,ref17,ref18}}，其中收缩因子的数值路径见文献\textsuperscript{\cite{ref13}}。表面条件经坐标换算后为

\begin{equation}
  -\frac{k}{R}\partial_\xi T\big|_{\xi=1}=h\,[T_R-T_\infty(t)],\qquad -D\,\partial_\xi \tilde C\big|_{\xi=1}=k_m R\,[\tilde C_R-C_\infty(t)]
  \label{eq:5-13}
\end{equation}
两点值得说明。其一，守恒形式~\eqref{eq:5-12} 中热容位于时间导数之内，其与“热容乘温度导数”的非守恒写法之差可显式写出：

\begin{equation}
  \frac{\partial(\rho c_pT)}{\partial t}=\rho c_p\frac{\partial T}{\partial t}+T\,\frac{\partial(\rho c_p)}{\partial t}
  \label{eq:5-25}
\end{equation}
即两者恰好相差右端第二项（热容随时间的变化率与温度之积 $T\,\partial_t(\rho c_p)$），而该项由脱水与体积压缩共同驱动，正是几何压缩项进入方程的代数来源；等价地，在固定物理位置上观察时，能量方程的时间导数会多出一项由半径变化率贡献的对流式作用，这就是几何压缩项进入方程的来源。该项是“热容随脱水衰减、同时又乘上体积压缩”这一耦合的物理来源，也是问题四出现压缩温升的机理；若沿用非守恒形式，这一项会被丢掉，压缩温升也随之消失。其二，界面系数仍取沿含水率的积分平均，其解析式与数值采样方式与问题二、三一致，以保证三个问题在同一离散口径下可比。

对报告值另作一次独立量级校验，即解析两界夹逼。下界由表面控制（集总极限）给出 8.36 h，该界严格成立；上界由固定半径数值解给出 129.48 h，它是同一问题在不收缩情形下的严格数值上界；另以渐近时标（半径平方与特征扩散系数之比，291.9 h）作量级校验。本文报告值 50.6523 h 落在两界之内，且与两界的相对位置符合“收缩加速干燥”的机理预期。

空间离散仍用元体平衡，径向区间数取 160（为问题二的 2 倍，因移动域的表面列收敛较慢）。时间上仍为交替推进，并对热容项采用预估–校正，使交付解中的热容项严格取更新层值，其依据是不做校正时物性冻结造成的滞后偏差达 O(1 K)，已超出四位小数点的取值预算；事件判定与子步级二分（精度 10\textsuperscript{-3} s）与问题三相同，并设 240 h 的熔断上界。

\subsubsection{结果与分析}
按题面格式给出每隔 6 h、若干位置（含药材表面）的水分浓度，见表~\ref{tab:6}。



\settabw{10}
\begin{longtable}{>{\centering\arraybackslash}m{\dimexpr 0.233\tblw\relax}>{\centering\arraybackslash}m{\dimexpr 0.183\tblw\relax}>{\centering\arraybackslash}m{\dimexpr 0.183\tblw\relax}>{\centering\arraybackslash}m{\dimexpr 0.183\tblw\relax}>{\centering\arraybackslash}m{\dimexpr 0.218\tblw\relax}}
  \caption{计入收缩后药材的水分浓度随时间和位置的变化（单位：kg/kg；各列为到药材中心的距离，末列为药材表面，单位：cm）}\label{tab:6}\\
  \toprule
    时间/h & 0 & 0.5 & 1.0 & 药材表面 \\
  \midrule
  \endfirsthead
  \multicolumn{5}{r}{\footnotesize（续）}\\
  \toprule
    时间/h & 0 & 0.5 & 1.0 & 药材表面 \\
  \midrule
  \endhead
  \bottomrule
  \endfoot
  \bottomrule
  \endlastfoot
  6 & 1.6276 & 1.4621 & 0.9868 & 0.4173 \\
  12 & 0.7021 & 0.6248 & 0.3930 & 0.1598 \\
  18 & 0.3958 & 0.3576 & 0.2340 & 0.0869 \\
  24 & 0.2794 & 0.2561 & 0.1763 & 0.0665 \\
  30 & 0.2232 & 0.2066 & 0.1477 & 0.0591 \\
  36 & 0.1909 & 0.1779 & 0.1306 & 0.0558 \\
  42 & 0.1699 & 0.1592 & 0.1193 & 0.0540 \\
  48 & 0.1552 & 0.1460 & 0.1110 & 0.0529 \\
  \textbf{烘干结束时间} & \textbf{0.1500} & 0.1413 & 0.1081 & \textbf{0.0526} \\
\end{longtable}

注 1：因半径收缩，1.5 cm 与 2.0 cm 两处自烘干开始后即越出药材边界（由附件 2 复算得 6 h 时半径约为 1.374 cm 且半径单调不增），故表~\ref{tab:6} 不设该两列，末列改为药材表面。该列契约与题面表~\ref{tab:6} 的模板一致，也与结果文件的列一致。

注 2：行次按题面表~\ref{tab:6} 模板，自 6 h 起，故不列 0 h 行；烘干结束行对应达标时刻，此时半径约为 1.200 cm，该行已由结果文件逐格取回。

计入尺寸收缩后的烘干时长为

\begin{equation*}
  t_{\rm dry}=182348.109\,\mathrm s=50.6523\pm0.0249\,\mathrm h
\end{equation*}
空间离散对烘干时长的贡献仅 0.0002 h，可以忽略；同一网格序列上表面含水率列仅为 0.99，即一阶，这正是本文对表面列额外加密的直接依据。时间三档呈非单调，故取保守上界 ±0.0249 h；合成口径为两者取较大者。

由表~\ref{tab:6}，可归纳出两条结论。

其一，收缩加快干燥。计入收缩后的 50.6523 h 短于问题三固定半径的 57.53 h。其净效应可分解为物性变化带来的延长与收缩带来的缩短两项，二者量级相当而方向相反（物性换套使时长延长约 71.9 h，收缩使其缩短约 78.8 h，净缩短 6.9 h；两项均为对照实验口径，不得作为独立结论引用）。这意味着几何收缩是本题中不可忽略的一阶效应，若沿用固定半径将明显高估烘干时长。

其二，守恒形式的能量方程揭示压缩温升，且中心仍恒为最慢点。最高温度达到 54.63 ℃，高于环境温度上限；该升温的时间位置与量级与准稳态估计一致，并被绝热包络严格覆盖。这意味着该现象是移动边界模型的内秉效应，而非数值振荡；其物理来源是体积压缩对热容的影响，即式~\eqref{eq:5-12} 中热容位于时间导数内所产生的附加升温。同时表~\ref{tab:6} 中每一时刻中心含水率均最高，48 h 时中心 0.1552 kg/kg、表面 0.0529 kg/kg，这意味着问题三的判据等价性结论在移动边界情形下依然成立。

本问的网格、时间与取法对照见表~\ref{tab:5-6}。表中两行值得强调：界面取法错误时网格加密并不收敛，故取法的正确性优先于网格分辨率；而时长分解属对照实验口径，不应作为独立结论引用。



\settabw{4}
\begin{longtable}{>{\raggedright\arraybackslash}m{\dimexpr 0.410\tblw\relax}>{\raggedright\arraybackslash}m{\dimexpr 0.590\tblw\relax}}
  \caption{问题四：网格、时间与界面取法对照}\label{tab:5-6}\\
  \toprule
    口径 & 结果 \\
  \midrule
  \endfirsthead
  \multicolumn{2}{r}{\footnotesize（续）}\\
  \toprule
    口径 & 结果 \\
  \midrule
  \endhead
  \bottomrule
  \endfoot
  \bottomrule
  \endlastfoot
  径向区间数 & 160（问题二的 2 倍，因移动域表面列收敛较慢） \\
  表面含水率列观测阶 & 0.99，即一阶 \\
  空间离散对烘干时长的贡献 & 0.0002 h \\
  时间误差带 & 取保守上界 ±0.0249 h \\
  界面取法（积分平均／调和平均） & 取法错误可使烘干时长被高估 52\% 以上 \\
  时长净效应分解 & 物性变化延长 71.9 h、收缩缩短 78.8 h、净缩短 6.9 h \\
  解析两界（下界／上界） & 8.36 h ／ 129.48 h \\
  常物性退化对拍最大温差 & $2.4\times10^{-10}$ K \\
\end{longtable}

注：界面取法与网格的二维对照见图~\ref{fig:A-10}；时长分解两项均为对照实验口径，不得作为独立结论引用。

压缩温升的机理与绝热包络如图~\ref{fig:5-20} 所示（该图置于表~\ref{tab:5-6} 之后，以便正文能从其右侧绕排）。


% 图5-20 ＝ F-49（F 号）／图45（PDF 文件索引号）／压缩温升机理与绝热包络
\begin{wrapfigure}{r}{0.500\textwidth}
  \centering
  \includegraphics[width=\linewidth]{../../05\_成品图/图45-Q4压缩温升机理与绝热包络.pdf}
  \vspace{2pt}
  \captionof{figure}{压缩温升机理与绝热包络}\label{fig:5-20}
\end{wrapfigure}

该图给出最高温度的时间位置、准稳态估计与绝热包络；由图~\ref{fig:5-20} 可见，该升温与准稳态估计在量级与时相上一致，并被绝热包络严格覆盖。

\subsubsection{模型检验（要点）}
问题四的检验包括网格收敛与时间误差带、常物性固定半径情形的级数解退化对拍、三项几何退化自检、程序化断言、局部与全局灵敏度、守恒核算、潜热对照、判据敏感度、低含水率端外推不确定度与界面取法对照。问题四的全局灵敏度中，低含水率端截断因子为第一主导因子（敏感度指数 0.337），即本问结论对低含水率段经验式的取法最为敏感，与 5.3.2 声明的适用边界相互指向。证据与图表见第六章。

本问另有两处断言须声明：压缩温升使中心温度超过环境温度上限，故朴素极值原理在此不成立，断言相应改为温度不超压缩包络上界；且移动边界下交替推进的物性滞后为 $O(1\,\mathrm K)$ 量级的系统偏差，问题四与固定域问题的差异属能量方程形式差异（模型级），不是实现缺陷。

\subsection{四问结果的横向对照与讨论}
前四节分别给出各问的建模、求解与结果。本节把四问的核心结果与可靠边界集中对照，不重复各问的结果呈现。为免歧义，本节所称含水率均指干基含水率。

四问核心结果对照见表~\ref{tab:5-3}。



\settabw{10}
\begin{longtable}{>{\raggedright\arraybackslash}m{\dimexpr 0.141\tblw\relax}>{\centering\arraybackslash}m{\dimexpr 0.244\tblw\relax}>{\centering\arraybackslash}m{\dimexpr 0.244\tblw\relax}>{\centering\arraybackslash}m{\dimexpr 0.169\tblw\relax}>{\centering\arraybackslash}m{\dimexpr 0.202\tblw\relax}}
  \caption{四问核心结果对照}\label{tab:5-3-a}\\
  \toprule
    项 & 问题一 & 问题二 & 问题三 & 问题四 \\
  \midrule
  \endfirsthead
  \multicolumn{5}{r}{\footnotesize（续）}\\
  \toprule
    项 & 问题一 & 问题二 & 问题三 & 问题四 \\
  \midrule
  \endhead
  \bottomrule
  \endfoot
  \bottomrule
  \endlastfoot
  物性体系 & 附录 2（常物性） & 附录 3 & 附录 3 & 附录 4 \\
  几何 & 固定，半径 2 cm & 固定 & 固定 & 收缩，2.000→1.198 cm \\
  时长口径 & 固定 1800 s & 固定 3 h & 反演（事件驱动） & 反演（事件驱动） \\
  时长结果 & — & — & 57.53±0.08 h & 50.6523 ±0.0249 h \\
  中心量（末端） & 温度 33.5753 ℃；含水率 2.5500 kg/kg & 温度 49.8495 ℃；含水率 1.7662 kg/kg & 含水率 0.1500 kg/kg & 含水率 0.1500 kg/kg \\
  表面量（末端） & 温度 36.7855 ℃；含水率 1.5104 kg/kg & 温度 49.9664 ℃；含水率 1.0081 kg/kg & 含水率 0.0526 kg/kg & 含水率 0.0526 kg/kg；半径 1.200 cm \\
  最高温度 & 36.7855 ℃ & 49.9664 ℃ & — & 54.63 ℃ \\
\end{longtable}

由表~\ref{tab:5-3} 可见，四问在物理上是同一过程的四个切片，结果的量级与趋势互为印证；“温度快、水分慢”这一时间尺度结构，也解释了为什么预热阶段的结果很早就接近环境条件，而干燥时长却由另一个更慢的过程决定。

跨问一致性共三条。其一，问题三与问题二在前 3 h 的逐时刻结果逐位一致，说明“复用加终止机制”未引入额外近似，长时程结果的偏差只能来自外推口径而非求解路径。其二，令问题四的半径取常值并使用附录 3 物性，其生产核与问题二结果一致至相对差 3.7×10\textsuperscript{-8}，已属舍入噪声量级，说明移动边界是加装式而非替换式。其三，问题一与问题二在 1800 s 处的系统性差异（中心温度差 −1.39 ℃、表面含水率差 +0.138 kg/kg）来自两问的物性换套，属模型必然而非计算错误。需强调的是，四问的衔接是“同量级且趋势连续”，而非“数值无任何跳变”。

从问题三的 57.53 h 到问题四的 50.6523 h，净缩短约 6.9 h，由两个方向相反的效应叠加而成：物性体系由附录 3 换为附录 4 使时长延长约 71.9 h，半径收缩（扩散路径变短）使其缩短约 78.8 h。两者量级相当而方向相反，故“收缩加快干燥”不能直接从结果表读出，须给出分解才能判断该结论对物性口径的敏感性；本文据此在第七章说明该结论的适用边界。

把四问结果连起来可得一条机理链：扩散系数随含水率下降而骤减，干燥因此在后期自我减速，烘干时长主要由低含水率段决定；收缩缩短了扩散路径，从而显著提前达标，而几何压缩在守恒形式下表现为压缩温升。工程上的启示是应以最慢点（中心）排产，并把控制重点放在温度与物料特性上。

本文结果的可靠边界见表~\ref{tab:5-4}（须与结论同页出现）。



\settabw{6}
\begin{longtable}{>{\centering\arraybackslash}m{\dimexpr 0.326\tblw\relax}>{\centering\arraybackslash}m{\dimexpr 0.299\tblw\relax}>{\centering\arraybackslash}m{\dimexpr 0.375\tblw\relax}}
  \caption{结果的可靠边界}\label{tab:5-4-a}\\
  \toprule
    类别 & 量级 & 说明 \\
  \midrule
  \endfirsthead
  \multicolumn{3}{r}{\footnotesize（续）}\\
  \toprule
    类别 & 量级 & 说明 \\
  \midrule
  \endhead
  \bottomrule
  \endfoot
  \bottomrule
  \endlastfoot
  离散误差带（问题三） & ±0.08 h & 空间网格收敛指数估计与时间保守上界取较大者 \\
  离散误差带（问题四） & ±0.0249 h & 同上；空间项可忽略，仅 0.0002 h \\
  低含水率端经验式外推（问题三） & 贡献 43\%–66\% 的时长 & 属最大不确定性来源 \\
  低含水率端经验式外推（问题四） & 贡献 16.8\%–55.1\% 的时长 & 其概率区间的先验为主观设定，引用须声明 \\
  界面取法 & 取法错误可使时长高估 52\% 以上 & 本文取沿含水率的积分平均 \\
  潜热忽略 & 温度系统性偏差 0.07–0.10 ℃ & 问题四的时长仅 +0.018\% \\
\end{longtable}

口径纪律：上述两类不确定度分列报告、不得相加；报告值一律取“空间与时间误差带取较大者”这一合成口径。表中除前两类离散误差带外，其余各项均为模型设定或经验式取法带来的系统性影响：前者随计算资源增加而收敛，后者只能通过口径声明与灵敏度分析加以界定——两者分开报告，即是不把可收敛误差与不可收敛误差混为一谈。效应分解的两项（+71.9 h／−78.8 h）均为对照实验口径，不是相互独立的答案，其差值与净效应仅用于解释变快或变慢的原因，不得作为独立结论引用。

另按题面要求逐项核对交付物落位：表~\ref{tab:1}／表~\ref{tab:2} 见 5.1.4，表~\ref{tab:3}／表~\ref{tab:4} 见 5.2.4，表~\ref{tab:5} 及烘干结束行见 5.3.4，表~\ref{tab:6}（含药材表面）见 5.4.4，四个结果文件（四位小数）列入附录支撑材料，均已生成并逐格核对；问题三、问题四的核心答案（烘干时间）均已带离散误差带给出。

从上述结果可以引出三点工艺含义。其一是控温上限的余量：由于收缩与变物性共同产生压缩温升，物料内部最高温度（54.63 ℃）可高于环境温度上限，故工艺控温设定不宜直接按环境温度取上限，应留出该内秉温升的余量。其二是排程口径：收缩使烘干时长较固定半径情形缩短约 61\%，若按常半径假设估算会系统性高估时长，故实际排程应以含收缩的模型结果为准。其三是停机判据：低含水率端经验式外推贡献了 43\% 至 66\% 的时长，末端时段的不确定性最大；本文认为，此时以含水率阈值判据停机、而非以定时停机，可把结论的不确定性限制在可控范围内。值得注意的是，这三点均来自本文已登记的对照实验口径，不构成新的独立结论。
